\section{Criteria for successful performance}
\label{sec:evaluation:success-criteria}
% Writing prerequisite for Chapter 7: finalize the accuracy targets and their
% rationale, including the plotted +/-5% DNS-relative target band. Distinguish
% an engineering acceptance target from DNS uncertainty or reporting precision.

Section~\ref{sec:evaluation:benchmark-cases} identifies uncertainties and ambiguities in the DNS references used for this assessment. In comparing RANS predictions with these references, as in related assessments such as that of Xue et al.~\cite{xue2024assessment}, a quantitative accuracy target is needed. This work adopts a \(\pm5\%\) DNS-relative error band as the target for wall-resolved heat-flux predictions. This band is an engineering acceptance criterion and neither represents the uncertainty of the individual DNS references nor defines a statistical confidence interval. The reference limitations discussed above must therefore be considered when judging agreement within this band.
\par\medskip

The target is motivated by the relative simplicity of attached ZPG flat-plate boundary layers and the accuracy already achieved by analytical approaches. Hasan et al.~\cite[Figure~2]{manzoor2024estimating}, for example, report heat-transfer errors mostly within \(\pm8\%\), with a maximum of \(10.3\%\), for a dataset that includes several of the benchmarks investigated here. Wall-resolved RANS should therefore aim to improve on this level of accuracy. For wall-function calculations, \(\pm10\%\) is adopted as the minimum acceptable heat-flux accuracy: for these ZPG cases, larger errors would leave the additional computational cost of CFD without a corresponding gain in heat-flux accuracy over the analytical approach.
\par\medskip

Beyond error magnitude, engineering use also requires an assessment of how consistently a method behaves. A systematic bias is easier to account for than an error that varies strongly and non-monotonically with sampling height or mesh spacing. This distinction helps frame the three analyses that follow: wall-resolved RANS discrepancies in Section~\ref{sec:comparative:zpg:wall-resolved-baselines}, wall-function closure errors in Section~\ref{sec:comparative:zpg:cross-comparison}, and spatial-discretization errors in Section~\ref{subsec:comparative:zpg:yplus-sensitivity}. For example, the wall-resolved heat-flux discrepancy for a given case may act primarily as a bias, while the closure and spatial discretization may introduce more irregular variations as the sampling height and mesh spacing change. The following sections examine these behaviours alongside the quantitative accuracy targets.
